% Источники: exam/materials/моделирование.pdf, раздел 7.
\section{Простые и структурно-сложные динамические системы.}

\textbf{Простыми динамическими системами} называются системы, поведение
которых описывается системами линейных дифференциальных уравнений в форме
Коши с достаточно гладкими правыми частями. Предварительная подготовка
решений обычно сводится к проверке совпадения числа неизвестных и числа
уравнений в системе.

Структура современных моделей часто соответствует структуре изучаемого
объекта. В основе таких моделей лежит элемент (блок), со скрытой от внешнего
наблюдателя внутренней структурой. Глядя на такой блок со стороны, можно
увидеть только специальные переменные, называемые в общем случае
контактными.

\textbf{Структурно сложная модель} состоит из множества блоков,
взаимодействующих между собой через функциональные связи между видимыми
извне переменными. Структура такой системы обычно является иерархической.
Множество элементов системы может изменяться в процессе функционирования
системы. Как правило, элементы сложной системы характеризуются различными
физическими принципами действия, что важно на этапе построения модели.

Наличие связи между контактами означает, что значения переменных,
соответствующих контактам, в любой момент равны. В современных визуальных
пакетах встречаются связи двух видов:
\begin{enumerate}
  \item однонаправленные: соединяемые контакты делятся на приемник и
  источник, а также постулируется, что приемник не может влиять на источник;
  \item двунаправленные: соединяемые контакты равноправны.
\end{enumerate}

\clearpage
